\section{Introduction}
\label{sec:intro}

Minimum-edge-cut recursive bisection draws congressional district maps by
repeatedly splitting the state's census-tract adjacency graph into two
population-balanced halves, minimising the total shared boundary length at
each level.
The algorithm makes no reference to partisan data: it sees only TIGER
boundary lengths and census populations.
Yet on states with concentrated metropolitan areas it reliably produces
a specific, troubling pattern.

\paragraph{The caterpillar pathology.}
Consider Illinois with 17 congressional districts.
The standard bisection always splits $\lfloor 17/2 \rfloor : \lceil 17/2 \rceil = 8:9$ at
the first level (as near-equal as possible).
But if we relax this constraint and ask METIS to try the $1:16$ split instead,
the edge-cut is dramatically cheaper: Chicago's compact urban boundary is
roughly 130\,km, while any boundary that bisects Illinois into
two $8$- and $9$-district regions must cross the state's full geographic width,
costing several hundred kilometres.

Minimum edge-cut thus \emph{always} prefers to carve off the densest compact
city before attempting a genuine bisection.
Illinois splits $1:16$ (Chicago peel).
New York splits $1:25$ (New York City peel).
Texas splits $1:37$ (Houston or Dallas peel).
Applied recursively, each level peels off the next-densest metropolitan core
until the remaining rural population is divided equally.
The result is a \emph{caterpillar bisection tree}: a linear chain of
$1:(k-1)$, $1:(k-2)$, $\ldots$ splits, one per urban core.

\paragraph{Why this is legally dangerous.}
The caterpillar pathology is not an algorithm bug --- each individual split
is a valid minimum-edge-cut bisection.
But the sequence looks identical to deliberate partisan packing:
one compact high-turnout Democratic district is created at every recursion
level, and the surrounding suburban ring districts land near the 50/50
partisan tipping point.
Under \emph{Rucho v.\ Common Cause} (2019), this may be unreviewable in
federal court, but state courts applying their own anti-gerrymandering
provisions (\emph{League of Women Voters v.\ Commonwealth of Pennsylvania}, 2018;
\emph{Harper v.\ Hall}, 2022) would scrutinise such a pattern.
An algorithm that produces caterpillar maps without any partisan input is
politically vulnerable, even if mathematically valid.

\paragraph{The isoperimetric fix.}
The cause of the pathology is straightforward: the raw edge-cut $\mathrm{EC}(i:k-i)$
is not a fair comparison across ratios because a $1$-district boundary is
intrinsically shorter than a $k/2$-district boundary.
For a convex region, the minimum boundary length to separate fraction $f$
of the population scales as $\sqrt{\min(f,\,1-f) \cdot A}$ where $A$ is total
boundary length, by the isoperimetric inequality for convex sets.
Dividing the edge-cut by $\sqrt{\min(i,\,k-i)}$ corrects for this scaling,
producing a normalised metric that treats all ratios on equal footing:
\[
  \mathrm{EC}_{\text{norm}}(i:k-i) = \frac{\mathrm{EC}(i:k-i)}{\sqrt{\min(i,\,k-i)}}.
\]
Under this metric, $1:k-1$ wins only if the corresponding urban boundary is
genuinely more compact per unit of geographic territory than any bisection.
For cities with compact municipal boundaries (Milwaukee, Philadelphia, Nashville),
$1:k-1$ still wins --- and this is the correct outcome, because those cities
are genuinely geographically isolated.
For states with distributed population (North Carolina's Piedmont corridor,
Wisconsin's rural west), a near-equal split wins because the urban cores are
not isolated enough to dominate the isoperimetric comparison.

\paragraph{Our contribution.}
We introduce \textbf{GeoSection}, a ratio-optimal recursive bisection algorithm
with two innovations:

\begin{enumerate}
\item \textbf{Isoperimetric ratio scan} (\S\ref{sec:algorithm}):
  at every bisection level, try all feasible ratios $1:(k-1)$ through
  $\lfloor k/2\rfloor:\lceil k/2\rceil$ with $N$ seeds each, select the
  ratio minimising $\mathrm{EC}/\sqrt{\min(i,k-i)}$.

\item \textbf{Recursive re-orientation} (\S\ref{sec:phase2}):
  each subregion after a split recomputes its own PCA minor axis and
  applies a directional penalty $\lambda$ to discourage cuts parallel to
  the subregion's narrowest extent.
\end{enumerate}

We make three additional contributions:

\begin{enumerate}\setcounter{enumi}{2}
\item \textbf{50-state empirical sweep} (\S\ref{sec:evaluation}): natural ratio
  results for all 45 multi-district states, showing the partition between
  states where $1:k-1$ is confirmed genuine and states where the
  normalisation shifts the winner to a more equal split.

\item \textbf{NC and WI case studies} (\S\ref{sec:case-studies}): full
  recursive bisection trees and district-level partisan analysis for the two
  most-studied states.

\item \textbf{Comparison to standard bisection} (\S\ref{sec:comparison}):
  head-to-head partisan outcome comparison against the MEC baseline (B.7),
  showing that geographic sorting dominates under both algorithms.
\end{enumerate}

\paragraph{Companion paper.}
GeoSection uses $\mathtt{ncon}=1$ (population balance only).
B.9 (AreaSection) adds a second METIS constraint requiring each half to
receive approximately 50\,\% of the state's land area ($\mathtt{ncon}=2$).
AreaSection is a separate algorithm with a different legal theory and
different empirical footprint; the two papers share notation and data.

\paragraph{Paper organisation.}
Section~\ref{sec:related} reviews related work.
Section~\ref{sec:algorithm} defines the GeoSection algorithm formally.
Section~\ref{sec:evaluation} presents empirical results across all 50 states.
Section~\ref{sec:discussion} discusses legal implications and limitations.
Section~\ref{sec:conclusion} concludes.
